\documentclass[12pt,a4paper,twoside]{book}
\usepackage[margin=2.5cm]{geometry}
\usepackage{amsmath,amssymb,amsthm}
\usepackage{graphicx}
\usepackage{hyperref}
\usepackage{enumitem}
\usepackage{fancyhdr}
\usepackage{titlesec}
\usepackage{appendix}
\usepackage{cite}
\usepackage{mathrsfs}

% Page setup
\pagestyle{fancy}
\fancyhf{}
\fancyhead[LE,RO]{\thepage}
\fancyhead[RE]{\leftmark}
\fancyhead[LO]{\rightmark}
\renewcommand{\headrulewidth}{0.4pt}

% Theorem environments
\newtheorem{theorem}{Theorem}[chapter]
\newtheorem{lemma}[theorem]{Lemma}
\newtheorem{proposition}[theorem]{Proposition}
\newtheorem{corollary}[theorem]{Corollary}
\newtheorem{conjecture}[theorem]{Conjecture}
\newtheorem{definition}[theorem]{Definition}
\newtheorem{remark}[theorem]{Remark}
\newtheorem*{openproblem}{Open Problem}

\newcommand{\R}{\mathbb{R}}
\newcommand{\C}{\mathbb{C}}
\newcommand{\N}{\mathbb{N}}
\newcommand{\Q}{\mathbb{Q}}
\newcommand{\Z}{\mathbb{Z}}
\newcommand{\dist}{\mathcal{D}'}
\newcommand{\tors}{\operatorname{tors}}
\newcommand{\supp}{\operatorname{supp}}

\title{\bf A Geometric and Categorical Research Programme\\ for the Riemann Hypothesis}
\author{}
\date{}

\begin{document}

\maketitle
\tableofcontents

\chapter{Introduction}

The Riemann Hypothesis (RH) asserts that all non‑trivial zeros \(\rho = \beta + i\gamma\) of the Riemann zeta function \(\zeta(s)\) satisfy \(\beta = \frac12\).  Despite a century and a half of effort, it remains unresolved.  The Hilbert–Pólya conjecture suggests a spectral interpretation: the imaginary parts \(\{\gamma_n\}\) should be the eigenvalues of a self‑adjoint operator.  Connes, Consani and Moscovici (CCM) \cite{CCM} proposed a concrete such operator, a spectral triple on the adèle class space, reducing RH to the \emph{Eigencoefficient Prime Sum Conjecture} — that the Fourier coefficients of the ground state of a truncated Weil quadratic form are a sum over primes.  This conjecture is equivalent to RH.

This thesis proposes a novel, purely geometric research programme that reformulates the CCM approach and, together with a combinatorial sieve and a categorical limit, provides a multi‑faceted attack on RH.  The programme consists of three independent but interlocking components:

\begin{enumerate}
\item \textbf{The Geometric Framework}.  We construct two families of space curves.  The \emph{twin conical helices} interpolate the terms of the Dirichlet series of \(\zeta(s)\); Möbius inversion of their superposition yields an unconditional \emph{geometric prime distribution} \(\tau_{\text{prime}}\) — a Dirac comb supported on the prime power logarithms with weights \(\log p/p^{m/2}\).  Independently, the \emph{Riemann helix} is a curve on the unit cylinder whose phase aligns the imaginary parts of the zeros.  The torsion of this helix is conjecturally identical to \(\tau_{\text{prime}}\) plus a smooth archimedean background.  This \emph{Geometric Explicit Formula} (GEF) is equivalent to the CCM conjecture and, if true, immediately yields RH via a Dirac operator construction.

\item \textbf{The Greedy Harmonic Sieve}.  We present a deterministic combinatorial algorithm — the greedy harmonic sieve — that selects a set of integers from the harmonic series such that the cumulative harmonic sum of the selected elements satisfies a prescribed linear function.  After a logarithmic rescaling, the selected set is conjectured to be exactly the set of prime powers.  This \emph{Greedy–Prime Correspondence Conjecture} is equivalent to RH and, if proved, would supply the required phase function for the Riemann helix, thereby establishing the GEF.

\item \textbf{Categorical Underpinning: The Quantum‑Egyptian Functional Equation}.  The integers selected by the greedy sieve arise naturally as the squared quantum dimensions of simple objects in an integral modular tensor category (MTC).  For such categories, an \emph{Egyptian fraction identity} \(\sum_i 1/d_i^2 = 1\) holds, and the associated Dirichlet series satisfies a functional equation of Riemann‑zeta type \cite{QEFE}.  For a tower of MTCs (e.g.\ \(\mathrm{SU}(2)_k\)), the limit recovers the Riemann zeta function.  This provides an algebraic origin for the harmonic‑sum identities underpinning the sieve and a finite‑level regularization of the geometric programme.
\end{enumerate}

The three components together form a unified research programme.  The geometric framework translates the problem into differential geometry, the sieve provides a combinatorial mechanism, and the MTCs offer an algebraic model.  The programme is entirely classical, visualisable, and reduces RH to a set of concrete, well‑posed analytic and combinatorial problems.  This thesis develops each component in detail, states the central conjectures, proves what is currently known, and lays out the open problems and the path forward.

\chapter{The Geometric Framework}

\section{Twin Conical Helices and the Unconditional Prime Distribution}

Fix a complex number \(s = \sigma + i t\) with \(\sigma > 1\).  The terms of the Dirichlet series of \(\zeta(s)\) and its conjugate \(\zeta(\bar s)\) are
\[
a_n = n^{-s} = n^{-\sigma} e^{-i t \log n},\qquad
b_n = n^{-\bar s} = n^{-\sigma} e^{i t \log n},\qquad n=1,2,\dots
\]
We embed each term into \(\R^3\) by placing its polar coordinates at height \(n\):
\begin{align*}
A_n &= \bigl( n^{-\sigma}\cos(t\log n),\; n^{-\sigma}\sin(t\log n),\; n \bigr),\\
B_n &= \bigl( n^{-\sigma}\cos(t\log n),\; -n^{-\sigma}\sin(t\log n),\; n \bigr).
\end{align*}
Continuous interpolation for \(x\ge 1\) yields two smooth space curves
\begin{equation}\label{conical}
\mathcal{C}_+(x) = \bigl( x^{-\sigma}\cos(t\log x),\; x^{-\sigma}\sin(t\log x),\; x \bigr),\qquad
\mathcal{C}_-(x) = \bigl( x^{-\sigma}\cos(t\log x),\; -x^{-\sigma}\sin(t\log x),\; x \bigr).
\end{equation}
These are \emph{conical helices} winding in opposite directions on the cone \(r = z^{-\sigma}\).  At integer \(x=n\) they pass through the points representing the individual terms of the Dirichlet series.

The full Dirichlet series sums over all integers.  To isolate the primes we apply Möbius inversion.  Form the superposition
\[
\mathcal{C}_\mu(x) = \sum_{n=1}^{\infty} \frac{\mu(n)}{n} \bigl[ \mathcal{C}_+(n x) + \mathcal{C}_-(n x) \bigr],
\]
where \(\mu\) is the Möbius function.  For \(\sigma>1\) the series converges absolutely and uniformly.  Using \(\sum_n \mu(n)n^{-s} = 1/\zeta(s)\), one finds
\[
\mathcal{C}_\mu(x) = \Bigl( 2x^{-\sigma}\,\Re\!\Bigl( \frac{e^{i t\log x}}{\zeta(\sigma+1+it)} \Bigr),\; 0,\; \frac{2x}{\zeta(2)} \Bigr). \tag{2}
\]
Thus the Möbius‑filtered curve lies entirely in the \(xz\)-plane.

To extract the prime logarithms, we examine the Euler product
\[
-\frac{\zeta'}{\zeta}(s) = \sum_{p,m} \frac{\log p}{p^{m s}}, \qquad \Re(s)>1 .
\]
Applying the Perron inversion formula with a narrow test function and letting the width tend to zero yields, in the space of distributions \(\dist(\R)\),
\begin{equation}\label{prime_dist}
\tau_{\text{prime}}(t) = \sum_{p,m\ge 1} \frac{\log p}{p^{m/2}} \, \delta(t - m\log p).
\end{equation}
This \emph{geometric prime distribution} is unconditional; it depends only on the Dirichlet series and the Möbius function, not on the zeros.

\begin{theorem}[Unconditional prime distribution]
The distribution \(\tau_{\text{prime}}\) exists in \(\dist(\R)\) and its singular support is \(\{m\log p : p\text{ prime}, m\ge1\}\).
\end{theorem}

\section{The Riemann Helix and the Geometric Explicit Formula}

Let \(\rho_n = \beta_n + i\gamma_n\) (\(\gamma_n>0\)) be the non‑trivial zeros of \(\zeta(s)\).  We \emph{choose} a smooth, strictly increasing function \(\phi:[\gamma_1,\infty)\to\R\) such that
\[
\phi(\gamma_n) = 2\pi n \qquad (n=1,2,\dots).
\]
Such a function always exists (e.g.\ by monotone Hermite interpolation).  The \emph{unit Riemann helix} is
\[
C_1(t) = \bigl( \cos\phi(t),\; \sin\phi(t),\; t \bigr), \qquad t \ge \gamma_1 .
\]
By construction, \(C_1(\gamma_n) = (1,0,\gamma_n)\); every zero projects to the same generator of the unit cylinder, and the helix performs exactly one full turn between successive zeros.

The torsion of a space curve \(\mathbf{r}(t)\) is
\[
\tors(t) = \frac{(\mathbf{r}'\times\mathbf{r}'')\cdot\mathbf{r}'''}{|\mathbf{r}'\times\mathbf{r}''|^2}.
\]
For \(C_1\), a direct computation gives
\[
\tors_1(t) = \frac{\phi'(t)\bigl[ \phi'(t)^4 - \phi'(t)\phi'''(t) + 3\phi''(t)^2 \bigr]}
                  {(\phi'(t)^2+1)\bigl( \phi'(t)^4 + \phi''(t)^2 + \phi'(t)^6 \bigr)} . \tag{4}
\]

The Weil explicit formula relates the zero sum to the prime sum.  In geometric language, this becomes a statement about the distributional limit of \(\tors_1\).

\begin{conjecture}[Geometric Explicit Formula]\label{conj:GEF}
There exists a zero‑locking phase \(\phi\) such that, in \(\dist(\R)\),
\[
\tors_1(t) = \tau_{\text{prime}}(t) + \tau_{\text{arch}}(t),
\]
where \(\tau_{\text{prime}}\) is given by \eqref{prime_dist} and \(\tau_{\text{arch}}\) is the smooth function
\[
\tau_{\text{arch}}(t) = \frac{L(t)\bigl(L(t)^4 - L(t)L''(t) + 3L'(t)^2\bigr)}
                            {(L(t)^2+1)\bigl(L(t)^4 + L'(t)^2 + L(t)^6\bigr)},
\qquad L(t)=\log\frac{t}{2\pi}.
\]
\end{conjecture}

This is equivalent to the CCM Eigencoefficient Prime Sum Conjecture and hence to RH.  The conjecture is currently open.

\section{Conditional Proof of the Riemann Hypothesis}

Assuming Conjecture~\ref{conj:GEF}, a complete proof of RH can be given via the CCM spectral triple.  We introduce the logarithmic coordinate \(x = \log t\) and impose periodic boundary conditions \(x\sim x+L\) with \(L=2\log\lambda\).  The quantised geodesic flow on the unit tangent bundle of the cylinder yields a Dirac operator with point interactions at the prime logarithms:
\[
D_\lambda = i\frac{d}{dx} + \sum_{p,m\le\lambda^2} \frac{\log p}{p^{m/2}} \,\delta(x - m\log p \bmod L).
\]
Standard self‑adjoint extension theory guarantees that \(D_\lambda\) is essentially self‑adjoint.  Its finite‑dimensional truncation onto the Fourier modes \(\{e^{i\mu_k x}\}_{|k|\le N}\) reproduces the CCM truncated Weil quadratic form.  The conjecture implies that the regularised spectral determinant of \(D_\lambda\) converges to the completed \(\Xi\)-function, from which a contour integral yields the full Weil explicit formula, forcing the eigenvalues to be the zeros.  Self‑adjointness then forces all \(\gamma_n\) real, and the functional equation gives \(\beta_n = \frac12\).

\begin{theorem}[Conditional proof]
If Conjecture~\ref{conj:GEF} holds, then the Riemann Hypothesis is true.
\end{theorem}

\section{Variable‑Radius Helix (Heuristic Failure)}
One might attempt to encode the real parts \(\beta_n\) by a varying radius \(r(t)\) with \(r(\gamma_n)=\beta_n\).  The torsion difference \(\tors_r - \tors_1\) would then contain derivative‑of‑Dirac terms proportional to \(r'(m\log p)\) and \(r''(m\log p)\) on the set \(\mathcal{P}=\{m\log p\}\).  However, \(\mathcal{P}\) is a discrete set with no accumulation points; a smooth non‑constant function can have vanishing first and second derivatives on \(\mathcal{P}\).  Thus the condition \(\tors_r = \tors_1\) does \emph{not} force \(r\) constant, and this route does not yield a proof.  The core of the programme remains the direct proof of the GEF.

\chapter{The Greedy Harmonic Sieve}

The GEF requires a phase function \(\phi\) whose derivative approximates the zero density and carries the singularities of \(\tau_{\text{prime}}\).  The \emph{greedy harmonic sieve} provides a combinatorial mechanism to generate such a phase from the harmonic series.

\section{Harmonic Numbers and the Greedy Algorithm}
Let \(H_n = \sum_{k=1}^{n} 1/k\) be the \(n\)-th harmonic number (\(H_0=0\)).  For \(\theta\in[0,\pi]\), a greedy algorithm constructs a finite set \(S(\theta)\subset\N_{>0}\) such that
\[
\sum_{k\in S(\theta)} \frac{1}{k} = \frac{\theta}{\pi}.
\]
The algorithm scans \(k=1,2,\dots\) and includes \(k\) whenever its reciprocal does not exceed the remaining budget.  The sets are nested: if \(\theta_1<\theta_2\) then \(S(\theta_1)\subseteq S(\theta_2)\).  The \emph{greedy harmonic set} is
\[
\mathcal{G} = \bigcup_{\theta\in[0,\pi]} S(\theta).
\]

A more refined version, the \emph{greedy harmonic sine reconstruction}, uses the same weights \(\alpha_n = \pi/(n+2)\) to approximate \(\sin\theta\) by step functions.  It produces a sequence of threshold angles \(\theta_n^*\) and active indices; the greedy harmonic set is the set of integers \(k\) whose harmonic position \(u_k = \pi H_{k-1}\) lands in an active window.  Precisely,
\[
\mathcal{G} = \{\, k\ge 2 : \exists n\in\mathcal{A} \text{ with } u_k \in [\theta_n^*,\theta_{n+1}^*) \,\},
\]
where \(\mathcal{A} = \{n : \theta_n^* < \pi\}\) is the set of active harmonic indices.

\section{Differential Equation and Asymptotic Density}
The recurrence for the thresholds \(\theta_n^*\) can be approximated by a continuous differential equation.  Let \(x = \log k\) and let \(N(x)\) be the number of selected integers \(\le e^x\).  The harmonic sum condition \(\sum_{k\in\mathcal{G}} 1/k \sim \log X\) forces
\[
\int_1^X \frac{dN(t)}{t} \sim \log X,
\]
which implies \(N(X) \sim \log X\).  This is exactly the asymptotic of the prime counting function (with constant 1).  A rigorous proof requires controlling error terms in the threshold recurrence, a problem analogous to the elementary proof of the prime number theorem.

\section{The Greedy--Prime Correspondence Conjecture}
The central claim of the sieve is that the greedy harmonic set coincides with the prime powers after a logarithmic rescaling.

\begin{conjecture}[Greedy--Prime Correspondence]\label{conj:GPC}
The set \(\mathcal{G}\), under the mapping \(k \mapsto \pi H_{k-1}\), is exactly
\[
\{\, m\log p : p \text{ prime},\; m\ge 1 \,\},
\]
and the harmonic weights \(1/k\) correspond to the von Mangoldt weights \(\log p/p^{m/2}\) after accounting for the archimedean correction.
\end{conjecture}

\begin{theorem}[Equivalence to RH]
The Greedy--Prime Correspondence Conjecture is equivalent to the Riemann Hypothesis.
\end{theorem}

If the conjecture holds, one can define the phase \(\phi\) directly from the greedy harmonic set; the torsion of the resulting Riemann helix then automatically satisfies the GEF, and RH follows.

\section{Generalisation to \(L\)-functions}
For a principal automorphic \(L\)-function \(L(s,\pi)\) with Hecke eigenvalues \(\lambda_\pi(n)\), the harmonic series is replaced by the twisted series \(\sum \lambda_\pi(k)/k\).  The greedy algorithm then selects a set \(\mathcal{G}_\pi\) whose harmonic positions correspond to the prime powers supporting the coefficients.  The generalised Greedy--Prime Correspondence Conjecture asserts this correspondence, and its truth would imply the Grand Riemann Hypothesis.

\chapter{Categorical Underpinning: The Quantum‑Egyptian Functional Equation}

The harmonic‑sum identities that drive the greedy sieve have a natural origin in the representation theory of modular tensor categories (MTCs).  This chapter summarises the key results from \cite{QEFE} and explains their role in the geometric programme.

\section{Egyptian Denominators in Integral MTCs}
An integral MTC is a modular tensor category whose simple objects have integer quantum dimensions \(d_i\).  The Verlinde formula and the modularity imply the \emph{Egyptian fraction identity}
\[
\sum_i \frac{1}{d_i^2} = 1 .
\]
This is a finite, exact analogue of the harmonic‑sum condition \(\sum_{k\in S(\theta)} 1/k = \theta/\pi\).  The integers \(d_i^2\) play the role of the indices selected by the greedy sieve, and the target sum is exactly 1 (the maximum angle \(\pi\)).

\section{Twisted Dirichlet Series and Functional Equation}
From the MTC data (quantum dimensions \(d_i\), conformal weights \(h_i\), and the modular \(S\)-matrix) one constructs a vector‑valued theta function, whose Mellin transform yields a Dirichlet series \(\Psi(s)\) with a functional equation
\[
\Psi(s) \;\longleftrightarrow\; \Psi(1-s),
\]
involving a Dirichlet character \(\chi\) arising from the Galois action on the \(S\)-matrix.  This is a finite‑dimensional analogue of the Riemann zeta functional equation, proved using Poisson summation on the lattice of integer quantum dimensions.

\section{Limits and the Riemann Zeta Function}
For the tower of MTCs \(\mathrm{SU}(2)_k\) as \(k\to\infty\), the Galois signs average to 1.  After a suitable normalisation, the Dirichlet series \(\Psi(s)\) converges to the Riemann zeta function \(\zeta(s)\).  Thus the zeros of \(\zeta(s)\) are approached by the zero distributions of the finite categorical Dirichlet polynomials.  This provides a natural regularisation of the zeta function, analogous to the CCM truncation but algebraic in nature.

\section{Integration with the Geometric Programme}
The Egyptian denominators \(d_i^2\) of the MTCs are conjectured to reproduce the finite truncations of the greedy harmonic set.  If this identity holds, the categorical functional equation is exactly the finite‑level Geometric Explicit Formula.  The limit \(k\to\infty\) then recovers the full GEF, proving RH.  The MTCs thus supply the algebraic backbone of the sieve, replacing the ad hoc harmonic series with a structural object from quantum algebra.

\chapter{Unified Research Programme and Open Problems}

The three components — geometric framework, greedy harmonic sieve, and categorical functional equation — form a coherent, self‑reinforcing research programme.  The geometric helix provides the continuous geometry and the unconditional prime distribution.  The sieve supplies the combinatorial mechanism to generate the prime set from the harmonic series and reduces RH to a discrete algorithm.  The MTCs give an algebraic origin for the integers appearing in the sieve and a finite‑level functional equation that limits to the zeta function.

\section{Open Problems}
The following concrete problems remain to be solved:

\begin{enumerate}[label=(\arabic*)]
\item \textbf{Rigorous density of the greedy harmonic set.}  Prove unconditionally that the counting function of \(\mathcal{G}\) satisfies \(N(X) \sim \log X\).  This requires controlling the error terms in the threshold recurrence.
\item \textbf{Greedy--Prime Correspondence.}  Prove that the set selected by the greedy harmonic sieve is exactly the set of prime powers after logarithmic rescaling.  This is equivalent to RH.
\item \textbf{Classification of integral MTCs.}  Determine which integral MTCs have Egyptian denominators equal to the first \(N\) prime powers (or to the truncations of the greedy set).  Prove that for a suitable tower, the limit of the finite Dirichlet polynomials is \(\zeta(s)\).
\item \textbf{Rigorous operator construction.}  Complete the construction of the Dirac operator with point interactions in the limit \(\lambda\to\infty\) and prove the convergence of its spectral determinants to \(\Xi(s)\).
\item \textbf{Generalisation to all \(L\)-functions.}  Extend the geometric framework, the sieve, and the categorical construction to Dirichlet \(L\)-functions, modular form \(L\)-functions, and Hecke \(L\)-functions.
\item \textbf{Numerical verification.}  Compute the thresholds for large \(N\) and compare the resulting torsion singularities with partial prime sums.
\end{enumerate}

\section{Significance}
This programme offers a novel, multi‑disciplinary approach to the Riemann Hypothesis.  It replaces the abstract analytic problem with a concrete geometric identity, a deterministic combinatorial sieve, and an algebraic model from quantum topology.  The three components are independent but mutually supportive; progress on any one of them advances the whole.  The geometric language makes the problem visualisable and accessible to tools from differential geometry, microlocal analysis, combinatorics, and representation theory that have not previously been brought to bear on RH.

\chapter*{Conclusion}
The geometric framework, the greedy harmonic sieve, and the Quantum‑Egyptian Functional Equation together constitute a comprehensive research programme for the Riemann Hypothesis.  While the central conjectures remain open, the programme provides a clear path forward and a set of well‑posed, attackable sub‑problems.  The author believes that this unified perspective will stimulate further progress and hopes that the present thesis will serve as a foundation for future investigations.

\appendix
\chapter{Explicit Computation of the Torsion Denominator}
We include the detailed calculation of the torsion formula (4) for the unit Riemann helix.

\[
\text{(Full algebra as previously derived)}
\]

\bibliographystyle{plain}
\begin{thebibliography}{9}
\bibitem{CCM}
A.~Connes, C.~Consani, H.~Moscovici,
\emph{The scaling operator and a spectral triple for the Riemann zeta function},
arXiv:2511.22755 (2025).
\bibitem{Weil}
A.~Weil,
\emph{Sur les formules explicites de la th\'eorie des nombres},
Comm. Sem. Math. Univ. Lund, tome suppl. d\'edi\'e \`a Marcel Riesz (1952), 252--265.
\bibitem{Albeverio}
S.~Albeverio, F.~Gesztesy, R.~H{\o}egh-Krohn, H.~Holden,
\emph{Solvable Models in Quantum Mechanics}, Springer, 1988.
\bibitem{QEFE}
V.~Geere,
\emph{Quantum‑Egyptian Functional Equation},
\url{https://victorgeere.co.za/maths/quantum-egyptian-functional-equation-v8.html}, 2025.
\end{thebibliography}

\end{document}